\section{Рабочий проект}

% Раздел~4. Цель: показать конкретную реализацию системы, ключевые классы,
% контрольные примеры работы и результаты тестирования. Это уровень <<кода>>:
% разделы и подсекции опираются на реальные файлы и модули репозиториев
% ispgestorserver и ispgestoradmin.

\subsection{Структура серверной части}

\subsubsection{Каталоги исходного кода}

% TODO: 0.5 страницы. Описать структуру проекта Laravel:
% app/Console/Commands, app/Events, app/Http/Controllers,
% app/Http/Middleware, app/Http/Requests, app/Http/Resources,
% app/Jobs, app/Models, app/Services, app/Traits, app/Policies;
% routes/api.php, routes/channels.php, routes/console.php;
% database/migrations, database/seeders, database/factories;
% config/*, tests/Feature, tests/Unit.

\subsubsection{Модели Eloquent}

% TODO: 2 страницы. Таблица 4.1 -- основные классы моделей и их назначение.
% (приведенная таблица должна включать класс, таблица БД, ключевые связи,
%  бизнес-методы). Привести 1--2 примера выразительного кода (например,
% Client::isActive(), Wallet::deposit(), Invoice::markAsPaid()).

% Заготовка таблицы 4.1 -- Описание моделей системы
% Колонки: Класс | Таблица | Назначение | Ключевые методы

\subsubsection{Контроллеры}

% TODO: 1.5 страницы. Таблица 4.2 -- перечень контроллеров.
% Колонки: Контроллер | Маршрут (префикс) | Назначение.
% Привести пример типового контроллера (AdminInvoiceController) с разбором
% методов index/show/store/update/destroy/generateAuto/charge.

\subsubsection{Сервисы и фоновые задачи}

% TODO: 1.5 страницы.
% Сервисы (app/Services/):
% -- MikroTikService -- инкапсулирует низкоуровневые вызовы routeros-api-php
%   (connect, command, disconnect);
% -- MikroTikQueueSyncService -- алгоритм синхронизации queue/simple;
% -- AutoBillingService -- генерация счетов, оплата с кошелька;
% -- ClientSuspensionService -- блокировка/разблокировка абонента;
% -- AutomationSettingsService -- чтение настроек планировщика из БД.
% Задачи (app/Jobs/):
% -- GenerateMonthlyInvoices, ProcessClientSuspension,
%    ProcessAutoReactivation, SyncMikroTikQueues.
% Привести пример кода задачи (rough listing 15--20 строк) с обработчиком
% handle() и тестовым покрытием в tests/Feature/.

\subsubsection{Промежуточные обработчики и политики доступа}

% TODO: 1 страница. Middleware:
% -- auth:sanctum -- проверка Bearer-токена;
% -- permission:slug -- проверка разрешения (CheckPermission);
% -- super_admin -- проверка роли суперадминистратора
%    (EnsureEmployeeSuperAdmin);
% -- throttle -- ограничение частоты обращений.
% Policies (app/Policies/) -- проверки на уровне модели (например,
% ClientPolicy::update($user, $client) -- сотрудник может править только
% клиентов своей зоны, если такое ограничение задано).

\subsection{Структура клиентской части}

\subsubsection{Каталоги исходного кода}

% TODO: 0.5 страницы. SvelteKit-проект: src/routes, src/lib/api,
% src/lib/components, src/lib/stores, src/lib/utils, static/, build/.

\subsubsection{Маршруты и страницы}

% TODO: 1 страница. Таблица 4.3 -- перечень маршрутов SvelteKit.
% Колонки: Маршрут | Файл | Назначение | Требуемые права.

\subsubsection{Компоненты пользовательского интерфейса}

% TODO: 1 страница. Описание ключевых компонентов:
% -- Menu.svelte (боковое меню);
% -- Chart.svelte (рукописный SVG-график с интерполяцией Catmull-Rom);
% -- KpiCards.svelte (карточки показателей);
% -- ModalVerFactura.svelte (просмотр и печать счета в jsPDF);
% -- ChatWidget.svelte (виджет чата с Reverb-подпиской);
% -- FirewallTable.svelte (таблица правил с режимом редактирования).

\subsubsection{Хранилища состояния}

% TODO: 0.5 страницы. Описание stores:
% auth.svelte.ts (токен, права, метод can), app.svelte.ts (UI-состояние),
% usuarios.svelte.ts, mikrotik-firewall.svelte.ts.

\subsubsection{Слой обращения к программному интерфейсу}

% TODO: 0.5 страницы. Модули src/lib/api/. Помощник fetchJsonWithRetry с
% экспоненциальной задержкой. Конфигурация PUBLIC\_API\_BASE.
% Кэширование статистики через cookie (срок жизни 1 сутки).

\subsubsection{Подписка на события в режиме реального времени}

% TODO: 1 страница. Инициализация Laravel Echo + pusher-js. Подписка на каналы:
% private('client.' + currentClient.id) -- уведомления о статусе;
% private('ticket.' + activeTicket.id) -- сообщения чата;
% private('employees') -- системные уведомления для сотрудников.
% Обработчики событий: показ всплывающих уведомлений (svelte-sonner),
% обновление списка сообщений в чате, обновление KPI на дашборде.

\subsection{Тестирование разработанной системы}

\subsubsection{Стратегия тестирования}

% TODO: 0.5 страницы. Иерархия тестов:
% Unit-тесты (tests/Unit) -- модели и сервисы в изоляции от БД;
% Feature-тесты (tests/Feature) -- запросы к API через TestCase
%   с in-memory SQLite либо тестовой MySQL-БД;
% Vitest (frontend) -- тесты Svelte-компонентов и утилит.
% Используется PestPHP 4.1 (синтаксис в стиле it()/expect()).

\subsubsection{Модульное тестирование серверной части}

% TODO: 1.5 страницы. Привести пример теста (PHP/Pest, 20--30 строк),
% покрывающего:
% -- генерацию номера счета (Invoice::generateInvoiceNumber);
% -- проверку метода Wallet::makePayment с отрицательным балансом;
% -- сценарий блокировки/разблокировки (ClientSuspensionService).
% Привести таблицу 4.X -- покрытие тестами по подсистемам.

\begin{figure}[ht]
\centering
% TODO: листинг кода теста (Pest).
\fbox{\parbox{0.7\linewidth}{\centering\textit{(Здесь будет листинг модульного теста, например, проверки генерации номера счета SRI)}}}
\caption{Модульный тест генерации номера счета}
\label{fig:unit_test_invoice}
\end{figure}

\subsubsection{Интеграционное тестирование RouterOS API}

% TODO: 1 страница. Описание стенда: один маршрутизатор MikroTik
% RB951Ui-2HnD (или CHR в виртуальной машине), доступ по IP 192.168.88.1,
% RouterOS 7.x. Тест-команда php artisan mikrotik:test проверяет:
% (1) подключение по 8728/TCP, (2) выполнение /system/identity/print,
% (3) создание тестовой очереди /queue/simple/add name="test_..." и
% последующее /queue/simple/remove, (4) чтение /interface/print,
% (5) право на запись в /ip/firewall/filter (создание и откат правила-заглушки).

\subsubsection{Системное тестирование}

% TODO: 4--5 страниц с серией контрольных примеров (скриншоты системы).

На рисунке~\ref{fig:login_page} представлена страница входа в панель администрирования.

\begin{figure}[ht]
\centering
% TODO: Insertar aquí captura de la página de login del panel administrativo
% (URL https://...admin/login). Debe mostrar: campo Email con validación,
% campo Contraseña con toggle de visibilidad, botón "Iniciar sesión",
% panel derecho con badges (MikroTik nativo, KPIs en vivo, Multi-zona,
% Facturación Auto, Control Firewall, Auto-cut, Roles, Alerts,
% Real-time Support). Estilo: glassmorphic dark.
\fbox{\parbox{0.7\linewidth}{\centering\textit{(Здесь будет скриншот страницы входа в систему)}}}
\caption{Страница входа в панель администрирования}
\label{fig:login_page}
\end{figure}

На рисунке~\ref{fig:dashboard} представлена главная страница (Dashboard).

\begin{figure}[H]
\centering
\includegraphics[width=\linewidth,keepaspectratio]{Dashboard}
\caption{Панель аналитики системы (Dashboard)}
\label{fig:dashboard}
\end{figure}

На рисунке~\ref{fig:clients} представлен реестр абонентов.

\begin{figure}[H]
\centering
\includegraphics[width=\linewidth,keepaspectratio]{Clientes}
\caption{Реестр абонентов с детальной карточкой}
\label{fig:clients}
\end{figure}

На рисунке~\ref{fig:plans} представлен каталог тарифных планов.

\begin{figure}[H]
\centering
% TODO: Insertar aquí captura del catálogo de planes (/planes).
% Debe mostrar: tabla con columnas Nombre, Velocidad descarga/subida,
% Precio mensual, Estado (Activo/Inactivo), Featured (estrella),
% botones de acción Editar/Desactivar. Modal de creación/edición de plan
% con todos los campos (name, slug, description, download_speed,
% upload_speed, symmetric checkbox, monthly_price, billing_cycle,
% mikrotik_queue_name).
\fbox{\parbox{0.7\linewidth}{\centering\textit{(Здесь будет скриншот каталога тарифных планов)}}}
\caption{Каталог тарифных планов}
\label{fig:plans}
\end{figure}

На рисунке~\ref{fig:invoices} представлен список счетов на оплату.

\begin{figure}[H]
\centering
% TODO: Insertar aquí captura del módulo de facturación (/facturas).
% Debe mostrar: filtros por rango de fechas y estado, tabla con columnas
% Número (formato 001-001-NNNNNNNNN), Cliente, Fecha emisión, Vencimiento,
% Total, Estado (Pendiente/Pagado/Vencido), badge colorido.
% Si es posible, mostrar también el ModalVerFactura con el preview del PDF
% que se generaría con jsPDF.
\fbox{\parbox{0.7\linewidth}{\centering\textit{(Здесь будет скриншот модуля управления счетами)}}}
\caption{Модуль управления счетами}
\label{fig:invoices}
\end{figure}

На рисунке~\ref{fig:invoice_pdf} представлен пример сгенерированного счета в формате PDF.

\begin{figure}[H]
\centering
% TODO: Insertar aquí captura del PDF de una factura generada
% (descargar desde el ModalVerFactura, exportada vía jsPDF + jspdf-autotable).
% Debe mostrar: cabecera con datos del emisor (RUC, nombre comercial,
% dirección, snapshot fiscal), datos del cliente (ID, nombre, dirección),
% número de factura, fechas, tabla de líneas (concepto del plan + cantidad
% + precio unitario), subtotal, impuesto (15% IVA Ecuador), total,
% método de pago, status badge.
\fbox{\parbox{0.7\linewidth}{\centering\textit{(Здесь будет скриншот сгенерированного счета в формате PDF)}}}
\caption{Печатная форма счета, сгенерированная средствами jsPDF}
\label{fig:invoice_pdf}
\end{figure}

На рисунке~\ref{fig:mikrotik_firewall} представлена страница управления правилами брандмауэра MikroTik.

\begin{figure}[H]
\centering
\includegraphics[width=\linewidth,keepaspectratio]{CortaFuegos}
\caption{Управление правилами брандмауэра MikroTik}
\label{fig:mikrotik_firewall}
\end{figure}

На рисунке~\ref{fig:mikrotik_queues} представлена страница мониторинга простых очередей.

\begin{figure}[H]
\centering
% TODO: Insertar aquí captura de /mikrotik/colas. Debe mostrar:
% tabla con columnas Nombre de la queue (formato plan-N-cliente-M),
% Target IP, Max-limit upload/download, Tráfico actual (bytes/sec),
% Estado (verde si activa). Botón "Sincronizar ahora" en la cabecera,
% "Limpiar huérfanos" con confirmación.
\fbox{\parbox{0.7\linewidth}{\centering\textit{(Здесь будет скриншот мониторинга простых очередей MikroTik)}}}
\caption{Мониторинг простых очередей маршрутизатора}
\label{fig:mikrotik_queues}
\end{figure}

На рисунке~\ref{fig:chat} представлена страница технической поддержки.

\begin{figure}[H]
\centering
\includegraphics[width=\linewidth,keepaspectratio]{SoporteTecnico}
\caption{Подсистема технической поддержки с обменом сообщениями}
\label{fig:chat}
\end{figure}

На рисунке~\ref{fig:roles} представлена страница управления ролями сотрудников.

\begin{figure}[H]
\centering
% TODO: Insertar aquí captura de /usuarios/roles. Debe mostrar:
% lista de roles (super_admin, operador_clientes, administrador_financiero,
% administrador_red, observador), matriz de permisos en cuadrícula
% (filas = módulos: usuarios/clientes/planes/facturas/configuracion/
% soporte/mikrotik; columnas = acciones: ver/crear/editar/eliminar/gestionar),
% checkboxes por celda.
\fbox{\parbox{0.7\linewidth}{\centering\textit{(Здесь будет скриншот матрицы ролей и разрешений)}}}
\caption{Управление ролями сотрудников и разрешениями}
\label{fig:roles}
\end{figure}

На рисунке~\ref{fig:client_portal} представлен личный кабинет абонента.

\begin{figure}[H]
\centering
% TODO: Insertar aquí captura del portal del cliente.
% Debe mostrar: header con nombre del cliente y botón logout,
% sección "Mi plan actual" (nombre, velocidad, precio, fecha de próximo
% cobro), sección "Mi saldo" con balance del wallet y botón "Recargar",
% sección "Últimas facturas" con tabla (número, fecha, monto, estado,
% botón descargar PDF), sección "Soporte" con botón "Abrir chat con soporte".
\fbox{\parbox{0.7\linewidth}{\centering\textit{(Здесь будет скриншот личного кабинета абонента)}}}
\caption{Личный кабинет абонента}
\label{fig:client_portal}
\end{figure}

\subsubsection{Результаты тестирования}

% TODO: 1 страница. Сводная таблица 4.X -- результаты тестирования по
% сценариям (прецедент, ожидаемый результат, фактический результат,
% оценка). Замечания по найденным дефектам и их устранению.
